\chapter{Einstellung der Slow-Kreise}
\section{Niedervolt Spannungsversorgung}
\todo{wurde gemessen}
\todo{werte}
\todo{bilder von kaputten sachen}

\section{Slow Pulse}

\subsection*{Linke Seite}

\jafps{tikz/slow_pulse_kontrolle}{Aufbau zur Kontrolle der Pulse im Slow-Kreis}{0.75}

Nun werden die Pulse aus dem Szintillator mit dem Oszilloskop betrachtet. Der Aufbau entspricht dabei dem in \fig{tikz/slow_pulse_kontrolle} gezeigten.
In \fig{links/slow_pulse} sind einige so aufgenommene Pulse zu erkennen.

\jafps{links/slow_pulse}{Slow-Pulse vom Linken Szintillator, auf Kanal 1 direkt am Szintillator abgegriffen, Kanal 2 Hinter dem Hauptverstärker.}{0.75}

Man erkennt, dass das Signal Direkt aus dem Szintillator die Form einer Zerfallskurve hat, mit einer relativ langen Pulsdauer, in der Größenordnung von einigen $\unit{\micro\second}$
